\chapter{Cel projektu i wprowadzenie teoretyczne}
\label{cha:chapter001}

\section{Cel projektu}

Celem ćwiczenia jest zaprojektowanie i zaimplementowanie systemu klasyfikacji
statycznych gestów dłoni odpowiadających literom amerykańskiego języka migowego
(ASL --- \textit{American Sign Language}). System o nazwie \textbf{mnist-translator}
pracuje w czasie rzeczywistym, korzystając z obrazu z kamery internetowej,
i jest w stanie rozpoznać 24 litery alfabetu ASL (A--Y, z wyłączeniem liter J
i~Z, które wymagają ruchu dłoni).

Szczegółowe cele projektu obejmują:

\begin{itemize}
    \item zapoznanie się ze zbiorem danych Sign Language MNIST --- jego strukturą,
          liczebnością klas i formatem próbek,
    \item opanowanie potoku przetwarzania obrazu: od przechwycenia klatki wideo
          przez detekcję dłoni za pomocą biblioteki MediaPipe, aż po ekstrakcję
          wektora cech dla klasyfikatora,
    \item wytrenowanie i porównanie trzech modeli uczenia maszynowego:
          \textbf{SVC z jądrem RBF}, \textbf{Lasu Losowego} oraz
          \textbf{konwolucyjnej sieci neuronowej (CNN)},
    \item opracowanie modułowej aplikacji konsolowej umożliwiającej trenowanie,
          ewaluację i uruchamianie wnioskowania zarówno offline (z pliku CSV),
          jak i na żywo (z kamery),
    \item sformułowanie wniosków o możliwościach zastosowania klasyfikatorów
          pikselowych do rozpoznawania gestów w warunkach rzeczywistych.
\end{itemize}

\section{Wprowadzenie teoretyczne}

\subsection{Zbiór danych Sign Language MNIST}

Sign Language MNIST jest adaptacją oryginalnego zbioru MNIST cyfr~\cite{lecun1998}
do zadania rozpoznawania liter języka migowego. Każda próbka to obraz dłoni
o rozmiarze $28 \times 28$ pikseli w skali szarości, spłaszczony do wektora
784~cech liczbowych (wartości pikseli 0--255).

\begin{table}[H]
    \centering
    \caption{Statystyki zbioru danych Sign Language MNIST}
    \label{tab:dataset}
    \begin{tabular}{lrr}
        \toprule
        \textbf{Podzbiór} & \textbf{Liczba próbek} & \textbf{Liczba klas} \\
        \midrule
        Treningowy & 27\,455 & 24 \\
        Testowy    & 7\,172  & 24 \\
        \bottomrule
    \end{tabular}
\end{table}

Litery J (indeks 9) i Z (indeks 25) zostały wyłączone ze zbioru, ponieważ
ich znaki ASL wymagają ruchu i nie mogą być reprezentowane przez pojedynczy
obraz statyczny. Mapowanie etykiet na litery opisuje wzór:
\[
    \text{litera} = \texttt{chr}\!\left(\texttt{ord}(\text{'A'}) + \text{label}\right),
    \quad \text{label} \in \{0,1,\ldots,8,10,\ldots,24\}.
\]

\subsection{Metody klasyfikacji}

\subsubsection{Maszyna wektorów nośnych (SVM/SVC)}

Maszyna wektorów nośnych~\cite{cortes1995} szuka hiperpłaszczyzny maksymalnie
rozdzielającej klasy w przestrzeni cech. W przypadku nieliniowych granic decyzji
stosowane są jądra (\textit{kernels}), które implicitnie mapują dane do przestrzeni
wyższego wymiaru. Jądro RBF (\textit{Radial Basis Function}) definiowane jest jako:

\begin{equation}
    K(\mathbf{x}_i, \mathbf{x}_j) = \exp\!\left(-\gamma \|\mathbf{x}_i - \mathbf{x}_j\|^2\right),
\end{equation}

gdzie $\gamma > 0$ jest parametrem szerokości jądra. Dla danych w przestrzeni
784-wymiarowej SVC z jądrem RBF jest w stanie uchwycić złożone relacje między
pikselami bez jawnej inżynierii cech.

\subsubsection{Las Losowy (\textit{Random Forest})}

Las Losowy~\cite{breiman2001} jest metodą łączącą (\textit{ensemble}) drzewa
decyzyjne. Każde drzewo trenowane jest na losowej próbce danych
(\textit{bootstrap}) i korzysta z losowego podzbioru cech przy każdym podziale.
Klasyfikacja odbywa się przez głosowanie większościowe:

\begin{equation}
    \hat{y} = \arg\max_{c} \sum_{t=1}^{T} \mathbf{1}[h_t(\mathbf{x}) = c],
\end{equation}

gdzie $T$ to liczba drzew, a $h_t$ to $t$-te drzewo decyzyjne. Las Losowy jest
odporny na przeuczenie i nie wymaga normalizacji danych wejściowych.

\subsubsection{Konwolucyjna sieć neuronowa (CNN)}

Konwolucyjne sieci neuronowe~\cite{lecun1989} są szczególnie efektywne dla danych
obrazowych, ponieważ uczą przestrzennych filtrów bezpośrednio z pikseli.
Architektura zastosowana w projekcie składa się z dwóch bloków konwolucyjnych
(\texttt{Conv2D} $\to$ \texttt{MaxPool2D}), warstwy spłaszczającej (\texttt{Flatten})
i gęstej warstwy wyjściowej z funkcją \textit{softmax}.

CNN operuje na siatce pikseli $28 \times 28$, co pozwala jej wychwytywać
lokalne wzorce (krawędzie, narożniki, faktury) niedostępne dla klasyfikatorów
operujących na spłaszczonym wektorze cech.

\subsection{Detekcja dłoni --- MediaPipe}

MediaPipe~\cite{lugaresi2019} jest platformą Google do uczenia maszynowego
na urządzeniach końcowych. Moduł \textit{Gesture Recognizer} wykrywa dłoń
w klatce wideo i zwraca 21 punktów kluczowych (\textit{landmarks}) w przestrzeni
3D. Na ich podstawie wyznaczany jest prostokąt otaczający (\textit{bounding box}),
który jest przycinany i przeskalowywany do formatu $28 \times 28$ pikseli.

\subsection{Normalizacja cech}

Wartości pikseli z zakresu $[0, 255]$ są normalizowane do przedziału $[0, 1]$
przez proste dzielenie:
\[
    x' = \frac{x}{255{,}0}.
\]
Zapewnia to zgodność formatu z danymi treningowymi ze zbioru Sign Language MNIST.